\chapter{Revisão bibliográfica}
\label{sec_Revisao_bibliografica}

Este capítulo apresenta a revisão bibliográfica que fundamenta a presente dissertação, com ênfase nos trabalhos relacionados às operações de rendezvous e berthing, à dinâmica relativa entre espaçonaves, à modelagem cinemática e dinâmica de manipuladores robóticos espaciais e às estratégias de guiamento, controle e planejamento aplicadas a esse contexto.

A literatura foi organizada em eixos temáticos, de modo a estruturar a discussão dos principais avanços já reportados e das abordagens mais recorrentes em cada tema. Essa organização permite situar, ao final do capítulo, o escopo específico desta dissertação em relação aos trabalhos revisados, com destaque para a integração entre dinâmica orbital relativa, dinâmica de atitude da espaçonave perseguidora e dinâmica acoplada de um manipulador robótico espacial nas fases de aproximação e interação com o alvo.


\section{Cinemática e modelagem de manipuladores espaciais}

A modelagem cinemática de manipuladores espaciais constitui uma etapa central para descrever o movimento relativo entre os elos, determinar a posição e a orientação do efetuador final e estabelecer as relações entre o espaço das juntas e o espaço cartesiano \cite{fonseca2019kinematics}. No contexto espacial, essa modelagem adquire complexidade adicional, pois o manipulador está acoplado a uma plataforma que não pode ser tratada, em geral, como uma base fixa convencional. Em vez disso, a formulação precisa considerar a interação entre o manipulador e a espaçonave, bem como o uso consistente de referenciais inercial, orbital e fixo ao corpo \cite{fonseca2019kinematics,fonseca2016dynamics}.

Nesse cenário, abordagens clássicas de cinemática direta e inversa continuam desempenhando papel fundamental. A convenção de Denavit--Hartenberg permanece amplamente empregada para a definição sistemática dos referenciais e para a obtenção das transformações homogêneas entre os elos, permitindo calcular a pose do efetuador final a partir das variáveis articulares \cite{fonseca2019kinematics}. Quando a cinemática inversa não admite solução analítica em forma fechada, o problema pode ser formulado como uma otimização, em que se busca a configuração articular que minimiza o erro entre a pose desejada e a pose obtida pelo manipulador \cite{santos2022uncertainty,santos2023physics}.

Entretanto, em manipuladores do tipo espaçonave, a formulação cinemática não se limita à descrição geométrica da cadeia serial. Como a movimentação das juntas pode induzir deslocamentos e rotações da plataforma, torna-se necessário compatibilizar a cinemática do manipulador com os movimentos da espaçonave e com os referenciais relevantes ao problema orbital \cite{fonseca2019kinematics,fonseca2016dynamics}. Essa característica é especialmente importante em operações de proximidade, captura e serviço em órbita, nas quais erros de posicionamento e orientação podem comprometer o desempenho da missão \cite{arantes2010far,fonseca2016dynamics}.

Além das formulações analíticas tradicionais, estudos recentes têm explorado estratégias baseadas em aprendizado de máquina para apoiar a resolução da cinemática inversa e o planejamento de movimento. Redes neurais e métodos como máquinas de vetores de suporte vêm sendo empregados para estimar configurações articulares, quantificar incertezas e melhorar a convergência de procedimentos de otimização, inclusive em situações em que os efeitos dinâmicos da plataforma e do manipulador são incorporados ao problema \cite{santos2022uncertainty,santos2023physics,santos2022machine}. Esses resultados indicam uma tendência de combinar modelos cinemáticos clássicos com técnicas computacionais orientadas por dados, visando ampliar a autonomia e a robustez operacional de manipuladores espaciais.

De modo geral, a literatura apresenta contribuições consistentes para a cinemática de manipuladores espaciais, tanto sob a perspectiva geométrica tradicional quanto sob abordagens computacionais mais recentes. Ainda assim, observa-se que parte dos trabalhos concentra-se na determinação da pose e na resolução da cinemática inversa, enquanto outra parte busca incorporar, de forma mais explícita, os efeitos do acoplamento entre manipulador e plataforma. Essa distinção é relevante para a presente dissertação, pois a modelagem cinemática constitui a base para a formulação posterior da dinâmica acoplada do sistema durante operações orbitais de proximidade \cite{fonseca2019kinematics,fonseca2016dynamics,santos2022uncertainty}. 

\section{Planejamento de trajetória e autonomia}

O planejamento de trajetória em operações espaciais constitui uma etapa central para garantir que a espaçonave, ou o manipulador acoplado a ela, execute a missão de forma segura e compatível com as restrições do problema. Em cenários de serviço em órbita, inspeção, captura e berthing, esse planejamento não se restringe à definição geométrica de um caminho, mas passa a envolver simultaneamente restrições cinodinâmicas, limites de atuação, risco de colisão, incertezas associadas ao alvo e requisitos crescentes de autonomia \cite{santos2022machine,capolupo2017heuristic,fiori2024modeling}.

Entre as abordagens encontradas na literatura, destacam-se as formulações baseadas em controle ótimo e otimização multiobjetivo. No trabalho de \cite{santos2022machine}, o problema de planejamento automático de trajetória de uma espaçonave com manipulador em órbita é formulado de modo a considerar, de forma simultânea, o deslocamento da plataforma, a manipulabilidade do braço, o torque máximo, o posicionamento do efetuador final, a prevenção de colisão entre o manipulador e a espaçonave e a minimização dos esforços de atuação sob incerteza na região de berthing. Essa formulação evidencia que, em robótica espacial, o planejamento de trajetória está fortemente associado ao tratamento simultâneo de múltiplos objetivos de desempenho e de restrições físicas do sistema \cite{santos2022machine}.

Outra classe importante de métodos corresponde às técnicas de motion planning baseadas em amostragem. Em \cite{capolupo2017heuristic}, algoritmos como Rapidly Exploring Random Trees e Sampling Based Model Predictive Optimization são empregados para construir sequências viáveis de manobras em operações de proximidade ao redor de pequenos corpos celestes, respeitando restrições cinodinâmicas e de trajetória. Os autores ressaltam que esse tipo de abordagem é particularmente útil em missões de maior nível de abstração, nas quais o problema não é definido apenas por waypoints previamente fixados, mas por objetivos mais amplos, como observação, mapeamento ou pouso autônomo em ambientes complexos \cite{capolupo2017heuristic}.

A incorporação de aprendizado de máquina ao problema também tem recebido atenção recente. Em \cite{santos2022machine}, uma estratégia de machine learning é empregada para apoiar a análise de cinemática inversa dentro do processo de planejamento ótimo, utilizando dados de treinamento e tarefas anteriores para melhorar a convergência do procedimento de otimização. De modo semelhante, \cite{santos2023physics} propõe uma abordagem physics informed machine learning em que uma Support Vector Machine é combinada a um modelo físico simplificado para estimar boas soluções iniciais para o problema inverso, incorporando efeitos dinâmicos do satélite e do manipulador ao processo de aprendizagem. Em ambos os casos, observa-se uma tendência de integrar modelos físicos e técnicas orientadas por dados para aumentar a eficiência computacional e apoiar níveis mais elevados de autonomia \cite{santos2022machine,santos2023physics}.

Além do planejamento propriamente dito, a literatura também apresenta contribuições relevantes no campo do guiamento autônomo sob restrições posicionais. O trabalho de \cite{fiori2024modeling} propõe um arcabouço de modelagem, simulação e controle para o movimento de uma pequena espaçonave nas proximidades de uma estação espacial, tratando o rendezvous com obstáculos físicos por meio de teoria de potenciais virtuais e de um formalismo geométrico baseado em grupos de Lie. Embora esse estudo não trate diretamente da presença de um manipulador robótico, ele contribui para a compreensão de estratégias de guiamento e controle aplicáveis a cenários de aproximação orbital em que a segurança da trajetória e a orientação da espaçonave são aspectos críticos \cite{fiori2024modeling}.

De modo geral, a literatura mostra que o planejamento de trajetória em missões espaciais vem incorporando, de forma crescente, múltiplos objetivos, restrições complexas e diferentes graus de autonomia. Ainda assim, muitos trabalhos tratam separadamente o planejamento do movimento do manipulador, o guiamento da plataforma ou a navegação em ambientes restritos, sem integrar de forma explícita a dinâmica orbital relativa, a atitude da espaçonave e o comportamento acoplado do manipulador em uma única formulação. Essa observação é relevante para a presente dissertação, pois reforça a necessidade de abordagens integradas para as fases de aproximação e interação em operações de berthing \cite{santos2022machine,capolupo2017heuristic,fiori2024modeling,santos2023physics}. 

\section{Encontro, aproximação e estimação de pose}

As operações de rendezvous em ambiente orbital envolvem desafios que ultrapassam o simples problema de deslocamento relativo entre duas espaçonaves. Em missões de serviço em órbita, inspeção, captura e berthing, a operação depende da articulação entre navegação relativa, definição de trajetórias seguras, estimação da pose do alvo e controle da aproximação final. Esses elementos tornam-se fortemente interdependentes, sobretudo em cenários de maior proximidade, nos quais erros de percepção, navegação ou controle podem comprometer a segurança e o sucesso da missão \cite{arantes2010far,fiori2024modeling}.

Entre os trabalhos relevantes nessa linha, destaca-se o estudo de Arantes et al., que analisa manobras far e proximity para uma constelação de satélites de serviço e incorpora, no mesmo contexto, estratégias de transferência orbital e estimação autônoma da pose de um satélite cliente não cooperativo \cite{arantes2010far}. Nesse trabalho, as manobras distantes são tratadas com base na formulação de Lambert, enquanto a fase de aproximação e trajetória final é analisada por meio das equações de Hill. Além disso, durante a fase de inspeção, a pose relativa do alvo é estimada por visão computacional monocular, utilizando informações geométricas previamente conhecidas do satélite cliente \cite{arantes2010far}.

Esse resultado evidencia que, em operações de rendezvous, o problema não se limita ao planejamento da trajetória relativa. Em alvos não cooperativos, a aproximação segura exige também a capacidade de reconstruir, com precisão adequada, a posição e a orientação relativas do alvo ao longo da missão. No trabalho de Arantes et al., a estimação de pose é formulada como um problema de minimização associado à correspondência entre características visuais do alvo e sua projeção no sistema de imagem, permitindo integrar percepção visual e navegação relativa em um mesmo arcabouço operacional \cite{arantes2010far}.

Além da vertente associada à percepção, a literatura também trata o rendezvous sob a ótica do guiamento e do controle em proximidade. O trabalho de Fiori et al. propõe um arcabouço de modelagem, simulação e controle para o movimento de uma pequena espaçonave nas proximidades de uma estação espacial, considerando restrições posicionais e a presença de obstáculos físicos \cite{fiori2024modeling}. A abordagem emprega potenciais virtuais e um formalismo geométrico baseado em grupos de Lie para descrever tanto o movimento roto-translacional quanto os campos de controle. Embora esse estudo não trate diretamente de manipuladores robóticos, ele contribui para a compreensão de estratégias de aproximação final em cenários nos quais a convergência para o alvo deve ser compatibilizada com restrições geométricas e operacionais \cite{fiori2024modeling}.

De forma mais ampla, a literatura mostra que o problema de rendezvous pode ser tratado em diferentes níveis. Em alguns trabalhos, a ênfase recai sobre as manobras orbitais e a dinâmica relativa; em outros, sobre a navegação local, a percepção do alvo ou o controle em regiões de maior sensibilidade operacional. Observa-se, assim, um avanço gradual na integração entre trajetória relativa, percepção e guiamento autônomo, embora muitos estudos ainda privilegiem apenas uma dessas dimensões \cite{arantes2010far,fiori2024modeling}.

Essa observação é relevante para a presente dissertação, pois reforça a necessidade de tratar de forma articulada a dinâmica relativa entre perseguidor e alvo, a fase de aproximação curta e os requisitos associados à interação final em operações de berthing. Em particular, interessa aqui a formulação de um problema em que a aproximação orbital, a atitude da espaçonave perseguidora e a presença do manipulador robótico sejam considerados de maneira consistente dentro do mesmo contexto de missão.

\section{Dinâmica e controle da espaçonave com manipulador robótico}

A dinâmica de uma espaçonave equipada com manipulador robótico constitui um tema central nas operações orbitais de proximidade, pois o movimento das juntas não altera apenas a configuração do braço, mas também afeta a atitude da plataforma e, em determinadas situações, o movimento do seu \gls{cdm}. Diferentemente de manipuladores terrestres, cuja base pode ser tratada como aproximadamente fixa para fins de modelagem, o manipulador espacial opera acoplado a uma plataforma livre, sujeita a reações dinâmicas decorrentes das forças e torques internos gerados pelo acionamento das juntas \cite{fonseca2016dynamics,fonseca2019kinematics}.

Nesse cenário, formulações lagrangianas têm sido amplamente empregadas para descrever a dinâmica do sistema acoplado. No trabalho de da Fonseca, Unfried e Lima, por exemplo, a modelagem de uma espaçonave do tipo manipulador durante a fase de aproximação de curta distância para rendezvous e docking/berthing é desenvolvida por meio da formulação lagrangiana com multiplicadores de Lagrange e função dissipativa de Rayleigh \cite{fonseca2016dynamics}. Essa abordagem permite obter não apenas as equações de movimento, mas também as forças e os torques de reação nas juntas, transmitidos ao longo da cadeia cinemática até a plataforma. Os autores mostram que, em microgravidade, esses efeitos perturbam tanto o movimento rotacional da espaçonave quanto a posição do seu \gls{cdm}, o que pode comprometer a precisão requerida nas operações de captura e berthing \cite{fonseca2016dynamics}.

Do ponto de vista do controle, a necessidade de estabilizar a atitude da plataforma enquanto o manipulador executa sua tarefa impõe exigências adicionais ao sistema de guiamento e controle. No estudo de da Fonseca, Unfried e Lima, as simulações computacionais combinam LQR para o movimento orbital linear relativo e PID para a estabilização da atitude durante a operação do manipulador \cite{fonseca2016dynamics}. Esse resultado evidencia que, em sistemas espaciais com manipuladores, o problema de controle não se restringe ao rastreamento de trajetórias articulares, mas envolve também a compensação dos efeitos reativos do braço sobre a plataforma e a preservação das condições necessárias à aproximação final \cite{fonseca2016dynamics}.

A coerência entre cinemática e dinâmica também desempenha papel importante nesse tipo de modelagem. O capítulo de da Fonseca, Pontuschka e Lima destaca que a formulação cinemática de manipuladores do tipo espaçonave deve incluir, além do referencial fixo ao corpo, os referenciais inercial e orbital, de modo a permitir a descrição consistente da atitude, da órbita e do movimento do manipulador dentro de um mesmo problema \cite{fonseca2019kinematics}. Essa observação é relevante porque a formulação dinâmica posterior depende diretamente da definição adequada desses referenciais e das transformações entre eles.

Além das abordagens centradas no sistema acoplado espaçonave--manipulador, a literatura também apresenta contribuições relevantes para o controle de aproximação orbital sob restrições geométricas e operacionais. O trabalho de Fiori et al. propõe uma formulação para modelagem, simulação e controle de rendezvous automatizado nas proximidades de uma estação espacial, considerando a presença de obstáculos e empregando potenciais virtuais e um formalismo geométrico em grupos de Lie \cite{fiori2024modeling}. Embora esse estudo não trate diretamente de um manipulador robótico acoplado, ele contribui para a compreensão de estratégias de guiamento e controle em proximidade orbital, especialmente em cenários nos quais a trajetória e a orientação da espaçonave precisam respeitar restrições posicionais durante a aproximação \cite{fiori2024modeling}.

De forma complementar, trabalhos recentes baseados em aprendizado de máquina têm buscado incorporar informações físicas do sistema a procedimentos de planejamento e resolução de problemas inversos. Os estudos de \cite{santos2022machine,santos2023physics,santos2022uncertainty} indicam que técnicas de machine learning e physics informed machine learning podem melhorar a convergência de algoritmos de otimização e apoiar o cálculo da cinemática inversa e do planejamento de trajetória em robótica espacial. Embora esses trabalhos não tenham como foco principal a formulação completa da dinâmica acoplada da espaçonave com o manipulador, eles sugerem uma tendência de integração entre modelagem física, controle e inteligência computacional em aplicações orbitais.

De modo geral, a literatura mostra que a dinâmica e o controle de espaçonaves com manipuladores robóticos vêm sendo tratados por diferentes perspectivas, que incluem formulações lagrangianas, técnicas clássicas de controle, estratégias de guiamento em proximidade e métodos computacionais orientados por dados. Ainda assim, permanece relevante o desenvolvimento de abordagens que articulem, de maneira consistente, a dinâmica orbital relativa, a atitude da espaçonave perseguidora e o comportamento acoplado do manipulador robótico durante fases críticas de aproximação e interação com o alvo. Essa necessidade está diretamente relacionada ao escopo da presente dissertação. 

\section{Métodos matemáticos e estratégias de controle}

A literatura relacionada às operações de rendezvous, berthing e robótica espacial emprega diferentes métodos matemáticos e estratégias de controle para representar o movimento relativo, descrever a dinâmica do sistema e conduzir manobras orbitais e tarefas de manipulação. Essa diversidade reflete o caráter multidisciplinar do problema, que envolve dinâmica orbital relativa, atitude, modelagem de manipuladores, controle, otimização e, em trabalhos mais recentes, técnicas computacionais orientadas por dados \cite{fonseca2016dynamics,arantes2010far,fiori2024modeling,santos2022machine,santos2023physics}.

Entre os métodos matemáticos mais recorrentes, destaca-se a formulação lagrangiana para a modelagem de sistemas acoplados. No trabalho de da Fonseca, Unfried e Lima, essa abordagem é empregada com multiplicadores de Lagrange e função dissipativa de Rayleigh para descrever a dinâmica de uma espaçonave do tipo manipulador durante a aproximação de curta distância, incluindo forças e torques de reação nas juntas e seus efeitos sobre a plataforma \cite{fonseca2016dynamics}. Essa formulação é particularmente relevante porque permite representar, de forma consistente, os acoplamentos entre os corpos do sistema e os efeitos dinâmicos decorrentes do movimento do manipulador \cite{fonseca2016dynamics}.

No tratamento do movimento relativo orbital, as equações de Hill--Clohessy--Wiltshire continuam desempenhando papel importante como ferramenta matemática para descrever a dinâmica linear relativa entre perseguidor e alvo em órbitas circulares ou quase circulares. No trabalho de \cite{arantes2010far}, essas equações compõem a base da análise das manobras far e proximity de satélites de serviço, em conjunto com estratégias de transferência orbital e aproximação ao satélite cliente \cite{arantes2010far}. De modo semelhante, o estudo de da Fonseca, Unfried e Lima também utiliza as equações de Hill como parte da formulação empregada nas simulações da aproximação final \cite{fonseca2016dynamics}.

Do ponto de vista do controle, observam-se tanto estratégias clássicas quanto formulações geométricas mais recentes. No estudo de da Fonseca, Unfried e Lima, técnicas LQR e PID são empregadas, respectivamente, para o movimento orbital linear relativo e para a estabilização da atitude durante a operação do manipulador \cite{fonseca2016dynamics}. Já no trabalho de Fiori et al., o rendezvous automatizado sob restrições posicionais é tratado com potenciais virtuais e um formalismo baseado em grupos de Lie, buscando lidar com obstáculos e com a orientação necessária à aproximação final \cite{fiori2024modeling}. Essas diferenças mostram que a escolha da estratégia de controle depende do tipo de missão, do modelo adotado e das restrições operacionais consideradas \cite{fiori2024modeling,fonseca2016dynamics}.

Além dessas abordagens, métodos de otimização e técnicas de aprendizado de máquina também têm sido incorporados como ferramentas auxiliares. Os trabalhos de \cite{santos2022machine,santos2023physics,santos2022uncertainty} indicam que essas técnicas podem apoiar o planejamento de trajetória e a resolução de problemas inversos, melhorando estimativas iniciais, convergência e eficiência computacional. Ainda assim, tais abordagens aparecem, em geral, como complemento às formulações físicas do sistema, e não como substituição da modelagem dinâmica e cinemática clássica \cite{santos2022machine,santos2023physics,santos2022uncertainty}.

De modo geral, observa-se que os métodos matemáticos e as estratégias de controle empregados na literatura abrangem desde formulações energéticas e equações lineares de movimento relativo até técnicas clássicas de controle, potenciais virtuais, grupos de Lie, otimização e aprendizado de máquina. Para a presente dissertação, essa revisão é relevante porque evidencia os principais instrumentos teóricos e computacionais disponíveis para sustentar a modelagem integrada da dinâmica relativa, da atitude da espaçonave perseguidora e do comportamento do manipulador robótico em operações de berthing.

\section{Contexto operacional de missões de rendezvous}

As missões de rendezvous em órbita ocupam posição relevante no desenvolvimento de tecnologias espaciais, tanto por sua importância histórica quanto por sua aplicação em operações contemporâneas de suporte, inspeção, manutenção, reabastecimento e serviço em órbita. Ao longo das últimas décadas, essas missões deixaram de ser tratadas apenas como demonstrações pontuais de capacidade orbital e passaram a integrar arquiteturas operacionais mais amplas, nas quais a aproximação segura entre veículos espaciais constitui etapa essencial para a execução da missão \cite{arantes2010far,technicalreportsoyuz}.

No contexto de missões tripuladas, o relatório técnico sobre a espaçonave Soyuz mostra que as operações de rendezvous envolvem diferentes escalas de problema, desde a inserção orbital e os ajustes de transferência até a navegação de proximidade e a fase final de encontro com a Estação Espacial Internacional. O estudo destaca a evolução de diferentes perfis operacionais, incluindo esquemas clássicos, rápidos e ultrarrápidos de rendezvous, bem como o uso das equações de Hill--Clohessy--Wiltshire nas operações de proximidade, evidenciando a busca por maior eficiência temporal, economia de recursos e robustez operacional \cite{technicalreportsoyuz}.

Além do contexto das missões tripuladas, a literatura sobre satélites de serviço mostra que o rendezvous também se tornou elemento central em missões não tripuladas voltadas à inspeção, observação, manutenção e atendimento de espaçonaves-alvo. No trabalho de \cite{arantes2010far}, por exemplo, são analisadas manobras far e proximity de uma constelação de satélites de serviço, combinando transferência orbital, dinâmica relativa e estimação autônoma da pose de um alvo não cooperativo por visão computacional. Essa abordagem ilustra uma mudança importante no contexto operacional do rendezvous, que deixa de se restringir ao encontro entre veículos cooperativos e passa a incluir cenários de serviço em órbita com maior grau de autonomia e maior complexidade operacional \cite{arantes2010far}.

Esse alargamento do contexto operacional também traz novos desafios de guiamento, controle e segurança durante a aproximação. Em operações modernas, não basta conduzir a espaçonave em direção ao alvo; é necessário respeitar restrições geométricas, corredores seguros, limitações dinâmicas e, em alguns casos, a presença de obstáculos na região de aproximação. O trabalho de \cite{fiori2024modeling} ilustra essa perspectiva ao tratar o rendezvous automatizado nas proximidades de uma estação espacial sob restrições posicionais e com obstáculos, considerando ainda fases distintas de aproximação no referencial \gls{lvlh}. Embora esse estudo não tenha foco em manipuladores robóticos, ele mostra que o contexto contemporâneo de rendezvous exige soluções integradas de modelagem, guiamento e controle em cenários mais restritivos \cite{fiori2024modeling}.

De modo geral, observa-se que o contexto operacional das missões de rendezvous evoluiu de operações predominantemente cooperativas e rigidamente planejadas para missões cada vez mais autônomas, flexíveis e orientadas a serviço em órbita. Essa evolução amplia a relevância de estudos voltados à dinâmica relativa, ao guiamento em proximidade, à percepção do alvo e à interação entre espaçonave e manipulador robótico. Para a presente dissertação, esse contexto é importante porque evidencia que a modelagem das fases de aproximação curta e interação final deve ser compreendida como parte de um problema operacional mais amplo, no qual precisão, segurança e autonomia se tornam requisitos centrais. 

\section{Lacunas da literatura e posicionamento da dissertação}

A revisão bibliográfica apresentada nas seções anteriores mostra que a literatura relacionada às operações de rendezvous, aproximação e robótica espacial já oferece contribuições relevantes em diferentes frentes, como modelagem cinemática de manipuladores espaciais, planejamento de trajetória, guiamento em proximidade, estimação de pose, formulações dinâmicas e estratégias de controle.
Entretanto, esses avanços aparecem, em grande parte, distribuídos em eixos específicos de investigação, de modo que muitos trabalhos concentram sua contribuição em apenas uma parcela do problema, como a cinemática do manipulador, o planejamento ótimo da trajetória, a navegação relativa, a estimação de pose do alvo ou a dinâmica acoplada entre braço e plataforma.

No campo da cinemática e da modelagem geométrica, observa-se uma base consolidada para a descrição do movimento de manipuladores espaciais, especialmente por meio da convenção de Denavit-Hartenberg, de transformações homogêneas e da formulação da cinemática direta e inversa \cite{fonseca2019kinematics}.
Também há avanços no emprego de técnicas de aprendizado de máquina para auxiliar a resolução da cinemática inversa sob incerteza ou para melhorar a convergência de procedimentos de otimização \cite{santos2022uncertainty,santos2022machine,santos2023physics}. Contudo, esses trabalhos tendem a concentrar-se na determinação da pose, na solução do problema inverso ou na eficiência computacional do processo, sem necessariamente tomar como foco principal a integração entre a dinâmica orbital relativa, a dinâmica de atitude da espaçonave e a resposta acoplada do manipulador robótico ao longo das fases críticas de aproximação.

No que se refere ao planejamento de trajetória e à autonomia, a literatura mostra abordagens capazes de lidar com múltiplos objetivos, restrições cinodinâmicas, prevenção de colisão e diferentes níveis de inteligência computacional \cite{santos2022machine,capolupo2017heuristic,fiori2024modeling}.
Ainda assim, parte expressiva dessas contribuições trata o problema sob a ótica do planejamento do movimento da plataforma, do manipulador ou do guiamento em cenários específicos, sem desenvolver de forma articulada uma formulação que contemple, simultaneamente, o movimento relativo entre perseguidor e alvo, a evolução da atitude da espaçonave e a atuação do manipulador em uma mesma estrutura dinâmica voltada à fase de aproximação curta e à interação final.

Em relação às operações de rendezvous e aproximação, há trabalhos relevantes que integram dinâmica relativa, manobras far e proximity e estimação autônoma da pose de alvos não cooperativos por visão computacional \cite{arantes2010far}. Esses estudos são importantes para mostrar que a navegação de proximidade depende não apenas do modelo orbital relativo, mas também da capacidade de percepção e de reconstrução da pose do alvo. Entretanto, a maior parte dessas abordagens não tem como objetivo principal modelar, de forma integrada, o efeito da presença de um manipulador robótico sobre a dinâmica da espaçonave perseguidora durante a aproximação final ou durante a fase de interação com o alvo.

Entre os trabalhos analisados, o estudo de da Fonseca, Unfried e Lima é o que mais se aproxima do escopo desta dissertação, ao tratar a dinâmica de uma espaçonave do tipo manipulador durante a aproximação de curta distância para rendezvous e docking/berthing, incluindo as forças e os torques de reação nas juntas e seus efeitos sobre a plataforma \cite{fonseca2016dynamics}. Essa contribuição constitui uma referência direta para o presente trabalho. Ainda assim, permanece relevante avançar na articulação entre dinâmica relativa orbital, atitude da espaçonave perseguidora e comportamento do manipulador robótico em uma formulação integrada voltada especificamente às operações de berthing.

Dessa forma, a principal lacuna identificada na literatura não está na ausência de estudos sobre os componentes do problema de forma isolada, mas na limitada articulação entre esses componentes em uma mesma estrutura de modelagem e simulação. Em particular, ainda são pouco frequentes formulações que representem, de maneira integrada, a dinâmica relativa entre perseguidor e alvo, a atitude da espaçonave perseguidora e a dinâmica acoplada do manipulador robótico durante as fases críticas de aproximação e interação final.

Nesse contexto, a presente dissertação se posiciona na interseção entre dinâmica orbital relativa, dinâmica de atitude e robótica espacial, com foco na modelagem integrada de uma espaçonave perseguidora equipada com manipulador robótico durante as fases de aproximação curta e operação final de berthing. Seu objetivo é conectar, em uma mesma estrutura de modelagem e simulação, elementos que frequentemente aparecem de forma separada na literatura, de modo a representar de forma mais consistente a influência do manipulador robótico sobre o comportamento dinâmico da espaçonave em operações orbitais de proximidade.

\begin{table}[H]
\centering
\caption{Posicionamento dos trabalhos revisados em relação aos temas centrais da dissertação.}
\label{tab_posicionamento_literatura}
\scriptsize
\renewcommand{\arraystretch}{1.3}
\begin{tabular}{|c|c|c|c|c|c|p{10cm}|}
\hline
\textbf{Ref.} & \textbf{A} & \textbf{B} & \textbf{C} & \textbf{D} & \textbf{E} & \textbf{Síntese} \\
\hline
1 & X &  &  &  &  & Base para modelagem cinemática, referenciais e transformações homogêneas. \\
\hline
2 & X &  &  &  & X & Cinemática inversa e redes neurais sob incerteza. \\
\hline
3 & X & X &  &  & X & Planejamento ótimo, manipulador espacial e eficiência computacional. \\
\hline
4 & X & X &  &  & X & Integra modelo físico simplificado e aprendizado de máquina. \\
\hline
5 &  & X & X &  & X & Guiamento heurístico e \textit{motion planning} em cenários espaciais complexos. \\
\hline
6 &  &  & X &  &  & Manobras far/proximity, dinâmica relativa e estimação autônoma de pose. \\
\hline
7 &  & X & X & X &  & Guiamento e controle em aproximação com restrições posicionais e obstáculos. \\
\hline
8 &  &  & X & X &  & Interação entre plataforma e manipulador em aproximação curta. \\
\hline
9 &  &  & X &  &  & Contexto operacional de rendezvous e dinâmica relativa em missões tripuladas. \\
\hline
\textbf{Diss.} & \textbf{X} &  & \textbf{X} & \textbf{X} &  & \textbf{Integra a dinâmica relativa orbital, a atitude da espaçonave perseguidora e a dinâmica acoplada do manipulador robótico no contexto das fases de aproximação curta e berthing.} \\
\hline
\end{tabular}
\end{table}

Na Tabela \ref{tab_posicionamento_literatura}, a coluna A corresponde à modelagem cinemática do manipulador espacial;
a coluna B refere-se ao planejamento de trajetória e às estratégias de guiamento;
a coluna C indica trabalhos voltados a rendezvous, aproximação e estimação de pose;
a coluna D reúne contribuições relacionadas à dinâmica acoplada e ao controle;
e a coluna E identifica abordagens que incorporam aprendizado de máquina.
As referências numeradas correspondem a: 1) \textit{Kinematics for Spacecraft-Type Robotic Manipulators}; 2) \textit{Uncertainty Quantification in Inverse Kinematics Computation of Space Robots by Neural Networks}; 3) \textit{A Machine Learning Strategy for Optimal Path Planning of Space Robotic Manipulator in On-Orbit Servicing}; 4) \textit{Physics Informed Machine Learning for Path Planning of Robots}; 5) \textit{Heuristic Guidance Techniques for the Exploration of Small Celestial Bodies}; 6) \textit{Far and Proximity Maneuvers of a Constellation of Service Satellites and Autonomous Pose Estimation of Customer Satellite Using Machine Vision}; 7) \textit{Modeling, Simulation and Control of a Spacecraft: Automated Rendezvous under Positional Constraints}; 8) \textit{Dynamics Analysis of a Manipulator-Like Spacecraft in the Close Range Approach for Rendezvous and Docking/Berthing}; e 9) \textit{Comprehensive Systems Analysis of the Soyuz Spacecraft}. 